\chapter{经典主题诗词文句汇编}

\section*{一、奋斗·理想}

\begin{enumerate}[leftmargin=*]
    \item 任何一个美好的愿望，想要变成现实，都需要我们付出不懈的努力和奋斗。
    
    奋斗是刘禹锡笔下“千淘万漉”的辛苦，
    
    奋斗是郑板桥笔下“咬定青山”的坚韧，
    
    奋斗更是陆游笔下“少壮工夫老始成”的一番耐心和决心。
    
    \item 燃，是人生的高峰，是理想的高光。
    
    “会当凌绝顶，一览众山小”，这是青春的奔马；
    
    “致君尧舜上，再使风俗淳”，这是思想的奔腾；
    
    “莫道桑榆晚，为霞尚满天”，这是夕阳的奔涌。
    
    \item 人生如风，生命如火。只要我们心怀梦想，眼里有光，就能点燃当下，照亮未来！
    
    正所谓，“自信人生二百年，会当水击三千里”！
    
    让我们一同以梦为马，共赴韶华！
\end{enumerate}

\newpage

\section*{二、家国天下}

\begin{enumerate}[leftmargin=*]
    \item 我们可爱的祖国，历史悠久，幅员辽阔。
    
    从“黄河落天走东海”到“青海长云暗雪山”，
    
    从“吹面不寒杨柳风”到“铁马秋风大散关”。
    
    有多少壮丽的山河就有多少豪迈的传奇；
    
    有多少伟大的人民就有多少光明的未来。
    
    踏遍青山人未老，风景这边独好。
    
    \item “大风起兮云飞扬，威加海内兮归故乡，安得猛士兮守四方！”
    
    什么是天下呢？
    
    对于蹒跚学步的孩子来说，父母的怀抱就是天下；
    
    对于求索未来的学子来说，知识的海洋就是天下；
    
    对于志在造福百姓的仁人志士来讲，江山和人民就是天下。
    
    “苟利国家生死以，岂因祸福避趋之”，
    
    天下者，百姓之天下，人民之天下。
    
    \item 天下是“迟日江山丽，春风花草香。”
    
    天下也是“大漠孤烟直，长河落日圆。”
    
    天下也是“江山如此多娇，引无数英雄竞折腰。”
    
    古人讲，天下兴亡，匹夫有责。
    
    天下是山河，是岁月，是英雄，也是你我。
    
    希望你我都能竹杖芒鞋，行走天下。
    
    问天下谁是英雄，数风流人物，还看今朝。
\end{enumerate}

\newpage

\section*{三、寻味人生}

\begin{enumerate}[leftmargin=*]
    \item “长江绕郭知鱼美，好竹连山觉笋香。”
    
    中国人对滋味的追寻，从美味开始，但又不止于美味。
    
    “日啖荔枝三百颗，不辞长作岭南人。”
    
    面对生活的磨难，诗人笑对人生，用美食去抚慰内心的苦涩。
    
    “闻说故园香稻熟，片帆归去就鲈鱼。”
    
    跨越山海，家乡的滋味，足以洗去诗人漂泊的风尘。
    
    一方灶火，点燃的是信念，照亮的是人生。
    
    \item 大千世界，千滋百味，美不胜收。
    
    比方说，有美酒的醇香——“兰陵美酒夜光杯，玉碗盛来琥珀光”；
    
    还有美食的鲜香——“夜雨剪春韭，新炊间黄粱”。
    
    饭虽然很普通，但是非常的甜。
    
    当然了，所有的人间美味，都是由于丰富的人生追寻和品味才会有的。有多么丰富的人生，就有多么丰富的品味。
    
    \item 我们每天都在品尝各种各样的味道：
    
    “烹羊宰牛且为乐，会须一饮三百杯”——那是最豪迈的味道；
    
    “莫笑农家腊酒浑，丰年留客足鸡豚”——那是最安心的味道；
    
    “西塞山前白鹭飞，桃花流水鳜鱼肥”——那是最风雅的味道；
    
    还有呢，“至乐无声唯孝悌，太羹有味是诗书”——那是最清淡绵长的味道。
    
    每种味道啊，都是从自然中采撷，经由生活去酝酿。
\end{enumerate}

\newpage

\section*{四、不忘初心}

\begin{enumerate}[leftmargin=*]
    \item 任何光明的未来，都离不开坚实的本来：
    
    “问渠那得清如许，为有源头活水来”；
    
    任何无畏的气概，都源自于坚定的本来：
    
    “千磨万击还坚劲，任尔东西南北风”；
    
    而任何卓越的成就，都离不开对本来的坚守。
    
    对知识思想的叩问，对大地的扎根，对梦想的坚持，都让我们从本来走向更辽远的未来。
    
    \item 凡有所本，必有所来。
    
    \item 遥想当年曹孟德已然暮年，但他依然壮心不已，因为“志在千里”是他的初心；
    
    陶渊明回归田园，非常地欣喜快慰：“户庭无尘杂，虚室有余闲”，因为“性本爱丘山”是他的本来。
\end{enumerate}

\newpage

\section*{五、远方}

\begin{enumerate}[leftmargin=*]
    \item “海上生明月，天涯共此时”——这是望月怀远，人间深情；
    
    “欲穷千里目，更上一层楼”——这是登高望远，人生壮志；
    
    “路漫漫其修远兮，吾将上下而求索”——这是爱国报国，任重道远；
    
    “红军不怕远征难，万水千山只等闲”——这是奋斗之路，行稳致远。
    
    面对远方，纵使崎岖坎坷，也坚信“柳暗花明又一村”；
    
    纵使风高浪急，也坚信“直挂云帆济沧海”。
\end{enumerate}

\newpage

\section*{六、情感}

\begin{enumerate}[leftmargin=*]
    \item 如果你问我，这世间最美好的感情是什么？
    
    我想刘禹锡在夔州写下的《竹枝词》就道出了怦然心动，人间美好。
    
    “去年今日此门中，人面桃花相映红”——那是惊鸿一瞥；
    
    “身无彩凤双飞翼，心有灵犀一点通”——那是一见倾心；
    
    “有美一人，见之不忘”——那是刻骨铭心。
    
    心动一时，电光石火；心动一世，海枯石烂。
    
    从一时走到一世，那是一生的相伴相依，是永生的相濡以沫。
\end{enumerate}

\newpage

\section*{七、追梦}

\begin{enumerate}[leftmargin=*]
    \item 飒，是声音的回响，是精神的激荡。
    
    它是金风与万物的共振，也是人生与天地的交响。
    
    “大鹏一日同风起，扶摇直上九万里”——这是平交王侯的豪气之飒。
    
    “休言女子非英物，夜夜龙泉壁上鸣”——这是巾帼红妆的英武之飒。
    
    “大江歌罢掉头东，邃密群科济世穷”——这是革命青年的担当之飒。
    
    飒，在世间每一个角落回荡，也激扬起每一位追梦者的理想风帆。
\end{enumerate}

\newpage

\section*{八、少年}

\begin{enumerate}[leftmargin=*]
    \item 自古沙场征战，英雄多出少年。
    
    从“弯弓辞汉月，插羽破天骄”的霍去病，
    
    到“马作的卢飞快，弓如霹雳弦惊”的辛弃疾，
    
    少年英雄，提笔扬鞭，锋芒所向，锐不可当。
    
    少年是一种精神，一种风采，更是一种情怀，一种境界。
    
    它是“天生我材必有用，千金散尽还复来”的自信；
    
    它是“会当凌绝顶，一览众山小”的胸怀；
    
    它是“指点江山，激扬文字”，是“粪土当年万户侯”的豪迈；
    
    它是“天地转，光阴迫，一万年太久，只争朝夕”的使命情怀。
\end{enumerate}



\section*{九、少年}

\begin{enumerate}[leftmargin=*]
    \item 这世界就是这么荒谬，只有尝过糖的人，才明白什么是苦。
    
    地层如同一座沙漏，计算着朝代的次序。没有“无足轻重”的数据，当一切归于尘土，谁能证明普通人曾经爱过。
    
    当一个伟大的时代来临的时候，其实有很多的小人物都在兢兢业业地做好着自己应该做的事。当我们每一个人都能把自己所承载的工作做到最好的时候，这个时代一定是个伟大的时代。
\end{enumerate}